\documentclass[ebook,10pt,oneside,openany]{memoir}
% \usepackage[utf8x]{inputenc}
% \usepackage[english]{babel}
\usepackage[brazilian]{babel}
\usepackage[utf8]{inputenc}
\usepackage[T1]{fontenc}
\usepackage{url}
\usepackage{titlesec}
\usepackage{amsmath}
\usepackage{amssymb}
\usepackage{diffcoeff}
\usepackage{graphicx}
\usepackage{float}

% \diffdef { pvrule } { op-symbol = \partial }
% for placeholder text
\usepackage{lipsum}

\usepackage{algorithm}
\usepackage{algorithmic}


% \usepackage{algpseudocode}
% \usepackage[⟨options⟩]{fancyhdr}
\usepackage{listings}

\renewcommand{\algorithmicrequire}{\textbf{Input:}}
\renewcommand{\algorithmicensure}{\textbf{Output:}}

\renewcommand{\algorithmicendif}{\textbf{end}}



\title{Anotações das Aulas de Ciência de Dados\\
  \large Fundamentos Matemáticos, Teoria de Aprendizado  \\
  e Aplicações em Python}
\author{Diego Pinheiro, PhD \\ Rodrigo Monteiro, PhD}

\begin{document}
\maketitle

\newpage
\tableofcontents
\part{Introdução}
\chapter{Motivação}

\begin{itemize}
  \item Por que estudar fundamentos?
\end{itemize}


\part{Notações}
\chapter{Estruturas unidimensionais e multidimensionais}

Para conseguirmos caminhar juntos neste livro,  precisamos estabelecer algumas notações e definições básicas relacionadas a estruturas unidimensionais e multidimensionais, como vetores e matrizes.

\section{Vetores}

Os vetores $\mathbf{x}$ e $\mathbf{w}$ são representados por letras minúsculas em negrito em formato coluna.

\begin{equation}
\begin{aligned}
\mathbf{x} = 
\begin{bmatrix}
    x_0 \\
    x_1 \\
    \vdots \\
    x_n
\end{bmatrix}, \quad
\mathbf{w} = 
\begin{bmatrix}
    w_0 \\
    w_1 \\
    \vdots \\
    w_n
\end{bmatrix} 
\end{aligned}
\end{equation}
\section{Matrizes}
As matrizes representam indivíduos ou sujeitos nas linhas e variáveis nas colunas. Ou seja, coletarmos 4 informações de 10 indivíduos, teremos uma matriz com 10 linhas e 4 colunas. As as 4 informações de cada indivíduo forem representas em um vetor coluna, para construir a matriz de dados, podemos empilhar em cada linha as transpostas desses  vetores colunas.

\begin{equation}
\begin{aligned}
X = \begin{bmatrix}
    \mbox{---}  \mathbf{x}^T_1 \mbox{---}  \\
     \mbox{---}  \mathbf{x}^T_2 \mbox{---}  \\
    \vdots \\
     \mbox{---}  \mathbf{x}^T_n \mbox{---}  \\
\end{bmatrix} 
\end{aligned}
\end{equation}
\part{Fundamentos Matemáticos}
\input{Parts/MathematicalFundations/mathematicalFundations}

\bibliographystyle{plain} % We choose the "plain" reference style
\bibliography{refs} % Entries are in the refs.bib file

\end{document}